\section{What the agent proposed}
\label{sec:proposals}

\paragraph{The agent and its sandbox.}
The editor was a fresh Claude Opus~5.5 session, run through the API. It had read-only access to the exploration
split E and no access to the project's records or to S, H or T-city. Its tools could run the model, record
activations, apply candidate edits and run Python. Analysis was limited to closed-form operations such as means,
covariances, projections, SVD and single linear solves. Gradients, optimisers and iterative loss-based fitting were
not allowed.

\paragraph{What an edit could change.}
An edit could change at most four components, where a component is the attention output projection
(\texttt{o\_proj}), the MLP output projection (\texttt{down\_proj}) or one of the two layer norms of a given layer.
Within a component the edit was either a dense update of rank at most 4 or a sparse edit of whole rows or columns
(at most 512 per candidate). The query, key, value, gate and up projections, the embedding, the unembedding, the
final norm and the rows for ``unknown'' tokens were off limits. Before testing an edit, the agent had to name the
component, state what it believed the component computes, predict outcomes for each row type and predict activity
in the unedited model. Each submission also fixed a primary outcome (its \emph{signed endpoint}), a most important
component and a scalar variable $f$ that its account says the component computes.

\paragraph{The run.}
The run took 73 turns, 16 model-pass requests and 63 Python runs, and saved six edits, of which the agent submitted
three. It lasted about 9.7 hours from start to the last submission, and its API cost was \$3.58.
\Cref{tab:submissions} summarises the submissions. The hypotheses are the agent's claims, recorded before any
evaluation; they are not established explanations.

\begin{table}[t]
\centering\small
\caption{The agent's three submissions. The hypotheses are abridged from the stored submission records; quoted
phrases are verbatim. ``Weights changed'' counts matrix entries altered by the edit.}
\label{tab:submissions}
\begin{tabularx}{\linewidth}{@{}lL{0.25\linewidth}Y@{}}
\toprule
\# & Edit & The agent's hypothesis (a claim, not an established explanation) \\
\midrule
1 & Zero the \texttt{o\_proj} columns of heads L32h2 and L33h10. Repair; about 0.52M weights changed. &
The two heads write ``a nearly constant `unknown' vector'' at the answer position, the same in every population.
This ``acts as an evidence-independent default-abstention prior''. Other late heads write the evidence city, more
strongly when the evidence is clearer. False abstention occurs when the city signal is only moderately weak, so the
prior wins. Removing the prior should recover false abstentions but not change unanswerable rows, where the
``unknown'' margin is much larger. \\
\addlinespace
2 & Rank-1 update of the layer-35 \texttt{o\_proj}. Repair; about 4.19M weights changed, from two 2048-dimensional
vectors. &
Layer-35 attention ``carries an `an answer is available in the memory/note' signal'' at the answer position,
almost as strong on false abstentions as on correct answers and absent on unanswerable rows. The model does not turn
this signal into suppressing ``unknown''. The edit reads the signal and writes it toward city tokens and away from
``unknown''. Because it sits in the last layer, ``it cannot be undone by self-repair''. \\
\addlinespace
3 & Double the same two heads' \texttt{o\_proj} columns. Diagnostic; about 0.52M weights changed. &
The converse of edit 1. If the heads supply the abstention prior, strengthening them should turn correct answers
into ``unknown'' while the note still names the city. \\
\bottomrule
\end{tabularx}
\end{table}

\paragraph{The selected repair in detail.}
Edit 2 adds a rank-1 matrix to the layer-35 output projection $W_O \in \mathbb{R}^{2048 \times 2048}$:
\[
  W_O \;\leftarrow\; W_O + s \, c \, d\,(W_O^{\top} w)^{\top}.
\]
Here $w$ is a read direction found by one closed-form solve on E. It was chosen to be large on answerable rows and
small on unanswerable rows. The write direction $d$ is a unit vector that the agent built from the difference between the
mean city unembedding and the ``unknown'' unembedding. The constant $c = 20$ is the agent's scale, and $s$ is the dose chosen on S. For an attention output
$a$, the edit adds $s\,c\,(w^{\top} W_O a)\,d$ to the residual stream: it reads the projection of the original
layer output onto $w$ and writes that amount along $d$. The update has rank 1 and is stored as two vectors. Its
dense matrix still changes almost every entry of $W_O$, which is why we describe it as a low-rank update and not as
a sparse one.

\paragraph{The agent's predictions.}
For edit 1, the agent predicted recovery of about 85\% of false abstentions, at least 97\% of correct answers
retained, and an increase in unsupported answers on related-distractor rows from about 1\% to at most 4\%. For
edit 2, it predicted 75--80\% recovery, 92--97\% of correct answers retained, and unsupported answering on those rows
staying at or below 3\%. For edit 3, it predicted that 30--40\% of correct answers would become ``unknown''. These
predictions were recorded for later scoring. They did not influence selection.
